\section{Resource Wiring in Program.cs}
\label{sec:wiring}

\subsection{Canonical Structure}

The complete canonical file for the Game scenario is:

\begin{lstlisting}[language={[Sharp]C}]
// src/Game.AppHost/Program.cs
var builder = DistributedApplication.CreateBuilder(args);

var postgres = builder.AddPostgres("postgres")
    .WithDataVolume(isReadOnly: false)
    .WithPgAdmin()
    .WithLifetime(ContainerLifetime.Persistent);

var gameDb = postgres.AddDatabase("gamedb");

builder.AddProject<Projects.Game_Web>("webapp")
    .WithReference(gameDb)
    .WaitFor(gameDb)
    .WithExternalHttpEndpoints();

builder.Build().Run();
\end{lstlisting}

Each builder call in this file is load-bearing. Table~\ref{tab:builder-calls} documents the responsibility of every call, the failure mode produced by its omission, and the Aspire component that implements it.

\begin{table}[H]
\centering
\caption{Builder calls in AppHost \texttt{Program.cs} and the consequences of omitting each.}
\label{tab:builder-calls}
\small
\renewcommand{\arraystretch}{1.15}
\begin{tabularx}{\textwidth}{@{}>{\raggedright\arraybackslash}p{4.5cm} >{\raggedright\arraybackslash}X >{\raggedright\arraybackslash}p{3.5cm}@{}}
\toprule
\textbf{Call} & \textbf{Effect} & \textbf{Omission failure mode} \\
\midrule
\texttt{AddPostgres("postgres")} & Pulls \texttt{postgres} Docker image; assigns resource name used as DNS alias & No Postgres container; web app cannot connect \\
\addlinespace
\texttt{.WithDataVolume()} & Creates/reuses named Docker volume & Database state lost on container stop \\
\addlinespace
\texttt{.WithPgAdmin()} & Pulls \texttt{dpage/pgadmin4}; auto-registers the server connection & No browser DB UI \\
\addlinespace
\texttt{.WithLifetime(Persistent)} & Container survives AppHost exit; reused on next \texttt{dotnet run} & Fresh empty DB on every developer session \\
\addlinespace
\texttt{.AddDatabase("gamedb")} & Registers logical DB; name must match \texttt{AddNpgsqlDbContext<>} argument & \texttt{ConnectionStrings\_\_gamedb} not injected; EF cannot resolve context \\
\addlinespace
\texttt{.WithReference(gameDb)} & Injects \texttt{ConnectionStrings\_\_gamedb} env var into web app & App starts without connection string; runtime exception on first DB call \\
\addlinespace
\texttt{.WaitFor(gameDb)} & Delays web app launch until Postgres health check passes & Race condition; app may start before DB is ready \\
\addlinespace
\texttt{.WithExternalHttpEndpoints()} & Binds HTTP/HTTPS on \texttt{localhost} outside Aspire reverse proxy & Web app only reachable through Aspire dashboard proxy \\
\bottomrule
\end{tabularx}
\end{table}

\subsection{The Postgres Resource Name and DNS Aliasing}

The string argument to \texttt{AddPostgres} --- \texttt{"postgres"} in the example --- becomes the DNS alias under which the container is reachable from other resources within the Aspire-managed network. pgAdmin uses this alias as the host when auto-configuring its server connection. The alias is also the base for the \texttt{ConnectionStrings\_\_postgres} host entry injected into dependent projects. It should be kept lowercase and DNS-safe (letters, digits, hyphens).

\subsection{The Database Resource Name and Connection String Injection}

The string argument to \texttt{AddDatabase} --- \texttt{"gamedb"} --- determines the key under which the connection string is injected into the web project's environment. Aspire injects it as \texttt{ConnectionStrings\_\_gamedb} (using double underscore as the hierarchy separator, consistent with the .NET configuration system). The \texttt{Aspire.Npgsql.EntityFrameworkCore.PostgreSQL} package reads this key automatically when \texttt{AddNpgsqlDbContext\textless AppDbContext\textgreater("gamedb")} is called in the web project's \texttt{Program.cs}. The two strings must match exactly; a mismatch produces a runtime \texttt{InvalidOperationException} with a message indicating that no connection string named \texttt{"gamedb"} was found.

\subsection{Persistent Containers and Developer Workflow}

\texttt{ContainerLifetime.Persistent} changes the container lifecycle in two specific ways. First, the container is not stopped when the AppHost process exits (\texttt{Ctrl+C} or IDE stop). Second, on the next \texttt{dotnet run}, Aspire detects the existing container by name and attaches to it rather than creating a new one. This means that any data written in session $n$ --- EF migrations, seeded rows, saved game state --- is available in session $n+1$. The developer's workflow collapses to: run, work, stop, resume; no database re-seeding is required.

The trade-off is that the container must be manually pruned when a clean state is desired. The command \texttt{docker rm -f postgres} (using the resource name) achieves this. Alternatively, removing the named volume with \texttt{docker volume rm} achieves the same effect at the data layer without removing the container.

\subsection{WaitFor and EF Migration Ordering}

\texttt{WaitFor(gameDb)} ensures the web application process does not start until the Postgres container passes its TCP health check.
